\scanpage{434}
\mspage{346}

\uline{Proposition} 4.13.25. Soient $\unsure{M|\mathcal M}$, $\pi$, $\unsure{\Phi|\mathscr I|\mathscr D}$ comme dans 4.13.23, avec $I_{\alpha_\pi}\in\Phi$ ($\alpha_\pi$ étant le plus petit ordinal \unsure{de|du} cardinal $\unsure{\Phi|\pi}$). Alors l'ens $\Sigma_\pi^{\prime\Phi}$ des sous-\unsure{catégorie|catégories} de $M$ $\Phi$-$\pi$-accessibles dans $M$ est stable par intersection de familles de cardinal $\leq\pi$.$^{*}$

% slanted
\margin{\unsure{Si|si} $\pi$ n'est pas un cardinal limite, il faut dire : familles de cardinal $\deleted{\unsure{\leq|{}=}}\!<\pi$. cf. 4.13.10.3., où se \unsure{trouve} un énoncé relatif, légèrement plus général\unsure{;|.}}

\deleted{\unsure{Les}} (\inserted{NB.} \unsure{la|les} bijection $\Sigma_\pi^\Phi\simeq\Sigma_\pi^{\prime\Phi}$ de 4.13.24 \unsure{respecte|respectent} les relations d'ordre naturelles d'inclusion, donc transforme intersections en intersections)

Se voit comme 4.13.7 (qui est contenu dans l'énoncé actuel), en utilisant le lemme 4.13.10.

\uline{4.13.26} Soient $M$, $\pi$, $\Phi$ comme dans 4.13.23, et soit $\underline S$ une sous-catégorie \uline{str. pleine} de $M$, $f:\underline S\to M$ le foncteur d'inclusion. \deleted{On suppose déjà que $M$ satisfait la \unsure{1°|1\textsuperscript{re}} de la \unsure{partie|page} \unsure{qui|préc.} \unsure{loc. cit.} \unsure{i.e.|$\pi$}, $M$ \unsure{est|et}} La condition que $\underline S$ soit $\Phi$-$\pi$-accessible dans $M$ se décompose en :

1) $\pi$ est utile pour $M$ \uline{et} pour $\underline S$.

2) $M$ et $\underline S$ sont stables par $\varinjlim$ de type $\Phi$.

3) $f$ commute aux dites limites.

4) $f$ est $\pi$-accessible

5) $f(\unsure{S|\underline S}_\pi)\subset M_\pi$.

% archive sheet number: 433

